\chapter{Introduction}

In this chapter, we introduce the problem investigated in this work, its motivations, main objectives, expected contributions, associated research questions, and the general organization of the text.

\section{Problem Description}

Humans create mathematical models to understand complex systems and predict possible outcomes. This is what Johannes Kepler did when he analyzed data on the movement of known planets in the solar system and discovered that the orbits of the planets were elliptical, not circular. Kepler published his laws of planetary motion in 1609 and 1619, which describe the motion of the planets around the Sun. These laws were a major milestone in the development of celestial mechanics and the understanding of the solar system.

Understanding the relationship between variables and transforming that knowledge into mathematical expressions is a crucial tool of the scientific method. Symbolic Regression (SR) attempts to solve the problem of finding a mathematical expression that best fits a dataset by finding relationships between the variables. This enables the establishment of relationships between the model extracted from the data and the phenomenon being analyzed, providing a means to better understand the phenomenon. SR methods are useful for tasks such as scientific modeling, finding physical or chemical models that help to understand more complex systems, or in engineering problems, such as optimizing processes to reduce energy costs. They can also help improve the interpretability of other machine learning models.

Other methods such as Linear Regression (LR), Decision Tree (DT), and Neural Networks (NN) can also find models that fit well in a dataset, however, each has its disadvantages. An LR presents a simple model, but it cannot have the same accuracy as other methods for nonlinear problems. A DT, such as the CART algorithm, can achieve better results and even clearly present the generated model; however, its approach requires the discretization of the problem, which creates abrupt "jumps" between value ranges. In cases with many dimensions, this can lead to an explosion of possibilities, making the generated tree difficult to understand. Models generated by NNs for regression problems can achieve excellent results, but in addition to being notoriously difficult to interpret, they also have a high computational cost.

In contrast to these methods, Symbolic Regression (SR) aims to produce an easily understandable model (a mathematical expression) while maintaining high accuracy. This makes it a valuable tool for experts in various fields. The problem of SR has been extensively studied in the literature and is considered a computationally challenging problem (NP-hard). This difficulty arises from the vast search space, which encompasses expressions of varying sizes and formats, along with numerous combinations of operators and operands. The issue becomes even more complex with the increase in the number of dimensions (independent variables) in the problem.

The most common approach to Symbolic Regression (SR) utilizes genetic algorithms, which efficiently explore the search space. Another promising approach is the use of neural networks to guide the discovery of equations, leveraging their learning capabilities and ability to represent complex patterns. Recently, models based on Large Language Models (LLMs) have been employed to find equations in an encoding and decoding process, where the dataset is the input and the mathematical expression that best fits the data is the output.

However, a challenge in this approach is adapting LLMs, trained to generate text in general, to the specific task of generating accurate mathematical expressions. While natural language can be ambiguous and subjective, mathematical expressions are computable, and their accuracy is easily measurable. In addition, the generated expression must strike a careful balance between fitting the observed data and maintaining simplicity. An overly complex expression risks overfitting—modeling noise rather than the underlying pattern—which impairs its ability to generalize to new data and reduces its human-readability.

In this work, we address these challenges by introducing a novel pipeline named Seriguela. The name is a playful nod to the pipeline's core components: Symbolic Regression, Language Models, and Reinforcement Learning Algorithms. Our approach adapts a large language model (specifically, GPT-2) to generate mathematical equations, which are then evaluated and refined using deep reinforcement learning. We systematically compare multiple RL algorithms---including REINFORCE, Proximal Policy Optimization (PPO), Group Relative Policy Optimization (GRPO), and hybrid Best-of-N methods---to identify the most effective approach for this domain. Furthermore, we investigate the robustness of these methods under challenging conditions, including noisy data and misleading prompt information. This comprehensive evaluation provides insights into both the capabilities and limitations of LLM-based symbolic regression.

\section{Objectives and Contributions}

The main objective of this work is to develop and evaluate a novel pipeline for symbolic regression that leverages the strengths of both large language models (LLMs) and deep reinforcement learning (DRL). Specifically, we aim to:

\begin{itemize}
\item Propose a novel pipeline for symbolic regression (SR) that integrates pre-trained LLMs with deep reinforcement learning algorithms.
\item Systematically compare multiple RL algorithms (REINFORCE, PPO, GRPO, and hybrid methods) for guiding expression generation.
\item Investigate the impact of different reward functions, penalty strategies, and temperature scheduling on model performance.
\item Evaluate the robustness of the approach under noisy data and adversarial prompt conditions.
\item Analyze the effect of model scale (Base, Medium, Large) on expression quality and validity.
\item Demonstrate the effectiveness of the proposed approach on the Nguyen benchmark suite.
\end{itemize}

\section{Research Questions}

In this research, we aim to answer the following research questions:

\begin{itemize}
\item \textbf{RQ1:} Which reinforcement learning algorithm (REINFORCE, PPO, GRPO, or hybrid methods) is most effective for guiding LLM-based symbolic regression?
\item \textbf{RQ2:} How do different reward functions (R² Clipped, Length-Penalized, SR-IC) affect the trade-off between expression accuracy and complexity?
\item \textbf{RQ3:} How robust is the LLM-RL approach to noisy data and misleading prompt information?
\item \textbf{RQ4:} How does model scale (124M to 774M parameters) influence expression validity and quality?
\item \textbf{RQ5:} Does the mathematical notation format (infix vs. prefix) significantly impact model performance?
\end{itemize}

\section{Organization of the Text}

Chapter 1 presents the description of the studied problem, motivations, objectives, contributions, and associated research questions. Chapter 2 discusses the main concepts and related work. Chapter 3 presents the materials and methods used in the project. In Chapter 4, the experiments performed are presented, and Chapter 5 concludes with the results and discussion of the experiment.